\documentclass[10pt]{article}
\usepackage[margin=0.85in]{geometry}
\usepackage{amsmath,amssymb,amsthm}
\newtheorem{theorem}{Theorem}
\title{A rank-one counterexample to conjecture 00000001952}
\author{earthking11}
\date{October 6, 2026}
\begin{document}
\maketitle
\section*{The original assertion}
The conjecture defines the minimum index of a finite-index subgroup of
$\operatorname{Out}(F_n)$ and claims its value is
$2^n\binom n2$. It gives no restriction $n\geq2$ or $n\geq3$.
We use the positive rank $n=1$, for which the asserted value is
$2\binom12=0$.

The group is not replaced by a numerical proxy. Set
$F_n$ equal to the free group on $n$ generators, let
$\operatorname{Inn}(F_n)$ be the image of the conjugation homomorphism
$F_n\to\operatorname{Aut}(F_n)$, and define
$\operatorname{Out}(F_n)=\operatorname{Aut}(F_n)/\operatorname{Inn}(F_n)$.
For any automorphism $\alpha$, conjugating the inner automorphism
$c_g:x\mapsto gxg^{-1}$ gives
$\alpha c_g\alpha^{-1}=c_{\alpha(g)}$, proving that the image is normal
and that this is the actual outer automorphism group.

\begin{theorem}
The proposed rank-one minimum-index value is false, both if all
finite-index subgroups are permitted and if only proper ones are permitted.
\end{theorem}
\begin{proof}
For any group $G$ and subgroup $H$, the left-coset space $G/H$ is nonempty:
it contains the identity coset. If it is finite, its cardinality
$[G:H]$ is therefore a positive integer.
If all finite-index subgroups are permitted, $G$ itself has index $1$.
Consequently the minimum finite subgroup index is exactly $1$, including
for the actual group $G=\operatorname{Out}(F_1)$, whereas the displayed
formula is $0$.

If ``subgroup'' is instead understood to mean ``proper subgroup'', any
attained finite minimum must still be positive because its witnessing
subgroup has a nonempty finite coset space. It cannot be $0$ either.
This second argument does not assert a particular proper minimum, or
assume a theorem identifying $\operatorname{Out}(F_1)$ with another group.
Thus neither reading can give the conjectured value.
\end{proof}

The extra claim about smaller indices arising as kernels cannot repair a
false numerical minimum. The counterexample is to the assertion as
written; it makes no claim about a modified problem restricted to higher
rank.

\section*{Formal verification and prior submission}
Lean uses Mathlib's \texttt{FreeGroup (Fin n)}, \texttt{MulAut}, actual
conjugation homomorphism, its image subgroup, normality, quotient group,
and subgroup index. It proves the equivalence of finite quotient and
nonzero index. Both minimum predicates require an actual witnessing
subgroup attaining the minimum, as well as minimality. The capstones
\texttt{conjecture\_01952\_false\_at\_one} and its \texttt{proper} version
negate the formula in these two conventions on the actual group.

Previous closed PR \#240 used a rank-three determinant-kernel argument,
but its Lean proof was only a numerical comparison and contained none of
the relevant groups, subgroups, or indices. The reviewer also noted that
the problem already mentions action kernels. This submission does not
reuse that determinant argument or its numeric certificate: it formally
constructs the actual outer automorphism group and refutes the rank-one
formula via genuine finite coset indices.
The project compiles on Lean 4.33.1 with pinned Mathlib, no unfinished
proofs, no native evaluation, and no additional axioms. The only printed
foundational dependencies are propositional extensionality, classical
choice, and quotient soundness. No brute-force computation is needed.
\end{document}
